\section{CLUTRR External Behavioral Validation}
\label{app:clutrr}

\subsection{Frozen Scope and Semantic Identity}

Both studies use CLUTRR Task~1 at relation length three, without distractor
facts.  A row contains three kinship-relation sentences and one queried
relation from 23 output classes.  Entity names are canonicalized by order of
appearance.  The semantic ID is
\[
(\text{target relation},
 \text{edge relation}_1,
 \text{edge relation}_2,
 \text{edge relation}_3),
\]
where the edges follow the ordered proof chain.  Sentence order and entity
surface names are excluded from this ID.

For each source story, all six sentence permutations receive one shared
\texttt{source\_group\_id}.  Eligible targets have at least 45 complete
groups.  The evaluation rows use the first 34 groups after deterministic
sorting: six permutations from groups 0--32 and two from group~33, for 200
rows per target.  Exposure uses the next 11 groups: six permutations from
groups 34--43 and four from group~44, for 64 rows per target.  Thus no source
group contributes to both exposure and evaluation.

Targets are paired only when they have the same output label.  For each seed,
a frozen SHA-256 bit assigns one member of every pair to each arm.  The two
arms share 22,720 common-base rows with no target semantic ID and append
1,280 exposure rows, giving 24,000 rows per run.  Label histograms, total
budget, tensor shape, initialization, epoch orders, and evaluation rows are
paired.

\subsection{Model, Training, and Primary Analysis}

The model is a byte-level Transformer with four layers, model dimension 128,
four attention heads, feed-forward dimension 512, dropout 0.1, pooling over a
classification token, and 23 output classes.  Training uses deterministic CPU
float32 execution, AdamW, batch size 128, learning rate
$3\times10^{-4}$, and 30 fixed epochs.  Each epoch contains 188 updates, so
the final checkpoint follows 5,640 updates.  There is no validation selection
or early stopping.

For seed $s$, let $\mathbf Z_s^A$ be the complete target-exposure assignment
in Arm~A and let $\mathbf Z_s^B=\mathbf 1-\mathbf Z_s^A$.  For target $f$,
$\mathbf Z_{s,f}^{(1)}$ is whichever arm assignment exposes $f$, and
$\mathbf Z_{s,f}^{(0)}$ is the complementary assignment.  Let
$A_{s,f}(\mathbf z)$ be row-weighted accuracy under complete curriculum
assignment $\mathbf z$.  The primary estimand is
\begin{equation}
\tau_{\rm pp}
=
\frac{100}{SF}
\sum_{s=1}^{S}\sum_{f=1}^{F}
\left[
A_{s,f}(\mathbf Z_{s,f}^{(1)})
-A_{s,f}(\mathbf Z_{s,f}^{(0)})
\right].
\end{equation}
Targets and seeds have equal weight; each target's 200 evaluation rows have
equal weight.  This target-indexed estimand compares two complete,
complementary curricula; it does not identify a one-target direct effect while
holding the other target assignments fixed.

The 20,000-replicate percentile bootstrap resamples model seeds, complete
same-label semantic-pattern pairs, and source groups within each selected
pair member.  Exposed and unexposed predictions use the same group draws, and
all rows from a sampled group occurrence remain together.  Percentiles use
NumPy's linear quantile method.  A separate analysis code path reconstructs
arm alignment from materialized assignments rather than trusting stored arm
labels.

The estimand itself is finite-population: it averages the eight frozen seeds
and 80 frozen targets.  For interval construction, the bootstrap treats the
observed seed labels, 40 complete pair labels, and source-group labels within
each target as exchangeable under their empirical distributions.  The
justification is procedural symmetry after freezing: each seed uses the same
architecture and training protocol, each selected pair receives the same
complementary intervention, and source-group IDs only identify permutations
of one source story.  Seed and pair resampling preserve their complete
clusters; group draws and duplicate seed occurrences are shared across arms.
The deterministic target-prefix and source-group allocation prevent a
design-based interpretation for all eligible CLUTRR patterns or stories.

The post-hoc sensitivity analysis uses the 640 seed-pattern effects
\(d_i\).  With residual
\(u_i=d_i-\bar d\), each cluster dimension \(c\) has CR1 variance
\[
\widehat V_c=
\frac{G_c}{G_c-1}\frac{1}{N^2}
\sum_{g=1}^{G_c}\left(\sum_{i\in g}u_i\right)^2.
\]
Clustering by model seed and complete pattern pair gives
\(\widehat V_{\rm two-way}
=\widehat V_{\rm seed}+\widehat V_{\rm pair}
-\widehat V_{\rm seed\times pair}\).
The resulting standard error is 2.6187 points.  A conservative two-sided
\(t_7\) interval is \([17.1469,29.5313]\) points, still above zero.  This
post-hoc sensitivity is not a new decision rule or a replacement for the registered
hierarchical bootstrap.

\subsection{V3 Pilot and Frozen Precision Decision}

V3 contains 20 pattern pairs, 40 targets, eight seeds, and two arms, for 16
runs and 128,000 saved evaluation predictions.  Before outcomes, the protocol
required any positive, negative, or equivalence conclusion to have a 95\%
interval no wider than two percentage points.

\begin{table}[H]
\caption{V3 primary result and registered decision.}
\label{tab:app-clutrr-v3}
\centering
\small
\begin{tabular}{lr}
\toprule
Quantity & Value \\
\midrule
Point estimate & $+15.3375$ pp \\
95\% percentile interval & $[+9.8769,+21.3693]$ pp \\
Interval width & $11.4924$ pp \\
Frozen maximum width & $2.0000$ pp \\
Registered outcome & positive but imprecise (precision criterion not met) \\
\bottomrule
\end{tabular}
\end{table}

The V3 estimand averages eight frozen seeds and 40 frozen targets.  Conditional
on that setting, its bootstrap treats those seeds, 20 complete pairs, and
within-target source groups as exchangeable empirical draws while preserving
clusters and paired arms.  Deterministic selection precludes design-based
generalization to all eligible CLUTRR patterns or stories.

The interval is directionally positive but fails the registered precision
criterion.  That decision is retained without reinterpretation.

After V3 was closed, a labeled exploratory analysis used bootstrap-Shapley
variance attribution and a seed-by-pair random-effects decomposition.
Pattern-pair resampling accounts for 94.40\% of bootstrap variance, compared
with 4.90\% for model seeds and 0.70\% for source groups.  The random-effects
estimates are $162.6402$ pp$^2$ for pair heterogeneity, $0.3490$ pp$^2$ for
seeds, and $48.6744$ pp$^2$ for the seed-by-pair interaction.

\subsection{Exploratory Effect Heterogeneity}
\label{tmlrs:app:clutrr-heterogeneity}

This exploratory (not preregistered) analysis resolves the V4 contrast into
its \TTwoPatterns{} semantic patterns.  The pattern effects have mean
\TTwoMean{} pp,
median \TTwoMedian{} pp, sample standard deviation \TTwoSD{} pp, and range
\TTwoMin{}--\TTwoMax{} pp; \TTwoPositive{}/\TTwoPatterns{} are positive.
Effect size is
strongly associated with unexposed accuracy
(Spearman $\rho=\TTwoSpearman{}$).  A linear regression on that baseline
alone has $R^2=\TTwoBaseRSquared{}$; adding target-relation indicators raises
it to $R^2=\TTwoRelationRSquared{}$.  The median normalized gain is
\TTwoNGMedian{}, so the raw spread is not only a scale artifact, although
available headroom explains much of it.

Mean effects are \TTwoInAbsentMean~pp when the target relation is absent
from the three edges (\TTwoInAbsentN{} patterns) and \TTwoInPresentMean~pp
when it is present (\TTwoInPresentN), and \TTwoDistinctOneMean,
\TTwoDistinctTwoMean, and \TTwoDistinctThreeMean~pp for one, two, and three
distinct edge types (\TTwoDistinctOneN, \TTwoDistinctTwoN, and
\TTwoDistinctThreeN{} patterns).  Neither covariate predicts the effect at
conventional levels (two-sided Spearman $p=\TTwoInP{}$ and
$p=\TTwoDistinctP{}$, respectively).
Target relation as a categorical grouping explains \TTwoEta{} of the effect
variance, but
leave-one-relation mean effects remain within
\TTwoLeaveMin{}--\TTwoLeaveMax{} pp.  Figure
\ref{tmlrs:fig:clutrr-heterogeneity} shows the full distribution.  Exposure
gains are largest where the unexposed model is far from ceiling, and
normalized gain is high for most patterns; the heterogeneity is thus mostly
a headroom effect rather than evidence for a particular compositional
mechanism.

A calibrated prospective projection finds that the original two-point width
would require about 634 independent pairs even with infinitely many seeds.
Only 139 eligible V3-disjoint pairs remain, and using all of them with
infinitely many seeds projects to a 4.2690-point width.  The strict-precision
replication route was therefore closed rather than consuming more seeds under
an infeasible criterion.

\subsection{Prospective V4 Design}

V4 treats V3 as a declared pilot and changes the scientific question from a
two-point precision target to directional confirmation.  The prospective
design effect is
\[
\min(0.5\times15.3375,\;9.8769)=7.66875\ \text{pp}.
\]
A 100,000-draw 64-bit permuted congruential generator (PCG64) Monte Carlo
simulation evaluates one to four 20-pair blocks.  The frozen selection rule
chooses the smallest block count for which estimated power and the 95\% Wilson
lower confidence bound are both at least 0.90.

For candidate pair count \(Q\in\{20,40,60,80\}\), set the frozen design
effect to \(\delta=7.66875\) points and draw
\[
\begin{aligned}
\widehat\tau_Q&=\delta+\mathrm{SE}_Q Z,\qquad Z\sim N(0,1),\\
\mathrm{SE}_Q
&=c\sqrt{
\frac{\sigma_s^2}{8}
+\frac{\sigma_p^2}{Q}
+\frac{\sigma_{sp}^2}{8Q}
}.
\end{aligned}
\]
The V3 random-effects inputs are
\(\sigma_s^2=0.3490014\),
\(\sigma_p^2=162.6401609\), and
\(\sigma_{sp}^2=48.6744361\) pp\(^2\), with frozen calibration
\(c=1.0067880\).  The seed effect is shared across pairs, the pair effect
across seeds, and the interaction is seed-by-pair.  Source groups are nested
within targets and shared across arms; their contribution is not a fourth
explicit random effect, but enters \(c\) through calibration to the full V3
hierarchical percentile width, which resampled source groups.  A simulated
experiment succeeds exactly when
\(\widehat\tau_Q\geq1\) and
\(\widehat\tau_Q-1.959964\,\mathrm{SE}_Q>0\).
Power is the fraction of 100,000 successful draws; the block-size rule additionally
requires its binomial 95\% Wilson lower endpoint to be at least 0.90.

\begin{table}[H]
\caption{V4 power analysis for the first two candidate sizes.}
\label{tab:app-clutrr-power}
\centering
\small
\begin{tabular}{rrrrr}
\toprule
Blocks & Pairs & Runs & Power & Wilson lower \\
\midrule
1 & 20 & 16 & 0.74397 & 0.74126 \\
2 & 40 & 32 & 0.95828 & 0.95702 \\
\bottomrule
\end{tabular}
\end{table}

The minimum-passing design therefore has two blocks, 40 pairs, 80 patterns,
and 32 runs.  All target semantic IDs are disjoint from V3.  The primary success criterion,
frozen before materialization and training, requires
\(\tau_{\rm pp}\geq1\) and a two-sided 95\% interval lower endpoint above
zero.

\subsection{V4 Execution Integrity and Result}

Materialization validates 80 targets, 16,000 evaluation rows across two
blocks, 5,120 exposure rows, the shared 22,720-row common base, and all 32
training manifests.  The frozen execution assets include a
$43{,}840\times185$ tensor cache, 32 arrays of 24,000 training indices,
two 8,000-row evaluation caches, eight initialization records, and 240
epoch-order records.

All 32 runs reach epoch 30.  The post-execution validator checks finite model
and optimizer state, 960 run-by-epoch order records, 256,000 predictions,
four-way RNG pairing across both blocks and arms within each seed, prediction
ordering, separate-code-path metric reconstruction, and artifact hashes.

\begin{table}[H]
\caption{V4 primary and descriptive results.  Sign counts do not define
additional hypothesis tests.}
\label{tab:app-clutrr-v4}
\centering
\small
\begin{tabular}{lr}
\toprule
Quantity & Value \\
\midrule
Point estimate & $+23.3391$ pp \\
95\% percentile interval & $[+17.7498,+29.2257]$ pp \\
Registered outcome & directional criterion met \\
Block effects & $+25.0109,\ +21.6672$ pp \\
Positive seeds & 8/8 \\
Positive pairs & 40/40 \\
Positive target patterns & 76/80 \\
Cross entropy, exposed minus unexposed & $-1.4938$ \\
Brier score, exposed minus unexposed & $-0.4224$ \\
ECE-15, exposed minus unexposed & $-0.2018$ \\
\bottomrule
\end{tabular}
\end{table}

The V4 point estimate exceeds one point and its interval lies above zero, so
the registered directional positive rule passes.  V3 and V4 are not pooled:
they use disjoint targets, answer different registered decision questions,
and retain separate bootstrap draws and classifications.

\subsection{Paper Artifact}

The compact CLUTRR bundle contains 51 checksum-managed files:
design and effect preregistrations, freeze validations, execution-integrity
reports, run-artifact hashes, exact bootstrap draws, pair/pattern/seed tables,
code snapshots, claim boundaries, source-file provenance, and a standalone
verifier.  It excludes local paths and large checkpoints.  Its
\texttt{SHA256SUMS} file has SHA-256
\texttt{210dddc02840a7bab7fe1215378d36ac\allowbreak86cc574a06aeca238482eab86ba3ef36}.
